\begin{figure}[htbp]
\centering
\begin{tikzpicture}[x=1cm,y=1cm,
  wire/.style={draftblue,thick},
  tr/.style={draw=draftblue,fill=white,minimum width=0.5cm,minimum height=0.36cm,font=\tiny},
  cell/.style={draw=draftblue,fill=draftlight,minimum width=1.05cm,minimum height=0.62cm,font=\tiny,align=center},
  blk/.style={draw=draftblue,fill=white,font=\scriptsize,align=center},
  lab/.style={font=\scriptsize}]
% ---------- (a) 6T cell ----------
\begin{scope}
\draw[wire] (0,0) -- (0,3.0) node[above,lab]{BL};
\draw[wire] (4,0) -- (4,3.0) node[above,lab]{$\overline{\mathrm{BL}}$};
\draw[wire] (-0.3,2.5) -- (4.3,2.5) node[right,lab]{WL};
% access transistors
\node[tr] (w1) at (0.75,1.4) {W1};
\node[tr] (w2) at (3.25,1.4) {W2};
\draw[wire] (0,1.4) -- (w1.west); \draw[wire] (w2.east) -- (4,1.4);
\draw[wire] (0.75,2.5) -- (w1.north); \draw[wire] (3.25,2.5) -- (w2.north);
% cross-coupled inverters
\node[draw=draftblue,fill=white,isosceles triangle,isosceles triangle apex angle=60,minimum size=0.5cm,inner sep=0pt] (i1) at (2,1.75) {};
\node[draw=draftblue,fill=white,isosceles triangle,isosceles triangle apex angle=60,minimum size=0.5cm,inner sep=0pt,shape border rotate=180] (i2) at (2,1.05) {};
\draw[draftblue,fill=white] (i1.apex) ++(0.06,0) circle (0.06);
\draw[draftblue,fill=white] (i2.apex) ++(-0.06,0) circle (0.06);
\draw[wire] (w1.east) -- (1.2,1.4);
\draw[wire] (1.2,1.4) |- (i1.west); \draw[wire] (1.2,1.4) |- ($(i2.apex)+(-0.12,0)$);
\draw[wire] (w2.west) -- (2.8,1.4);
\draw[wire] (2.8,1.4) |- ($(i1.apex)+(0.12,0)$); \draw[wire] (2.8,1.4) |- (i2.west);
\node[lab,align=center] at (2,0.35) {two cross-coupled\\inverters (4 transistors)};
\node[font=\small\bfseries,text=draftblue] at (2,3.95) {(a) 6-transistor SRAM cell};
\end{scope}
% ---------- (b) array ----------
\begin{scope}[xshift=6.3cm]
\node[blk,minimum width=1.0cm,minimum height=2.3cm] (dec) at (0.4,1.65) {Row\\de-\\coder};
\foreach \r in {0,1,2}{
  \foreach \c in {0,1,2,3}{ \node[cell] at (1.75+\c*1.2,2.5-\r*0.85) {cell}; }
  \draw[wire,draftblue!60] (0.9,2.5-\r*0.85+0.2) -- (6.15,2.5-\r*0.85+0.2);}
\node[lab,anchor=west] at (6.2,2.7) {word lines};
\foreach \c in {0,1,2,3}{\draw[draftblue!60] (1.3+\c*1.2,3.0) -- (1.3+\c*1.2,0.35); \draw[draftblue!60] (2.2+\c*1.2,3.0) -- (2.2+\c*1.2,0.35);}
\node[blk,minimum width=4.75cm,minimum height=0.42cm] at (3.55,0.05) {column multiplexer};
\node[blk,minimum width=4.75cm,minimum height=0.42cm] at (3.55,-0.5) {sense amplifiers and write drivers};
\draw[-{Stealth},thick,draftblue] (-0.5,1.65) -- (dec.west) node[pos=0,above,lab]{address};
\draw[-{Stealth},thick,draftblue] (3.55,-0.72) -- (3.55,-1.15) node[below,lab]{data};
\node[lab] at (3.55,3.25) {bit-line pairs (one per column)};
\node[font=\small\bfseries,text=draftblue] at (3.0,3.95) {(b) Array of SRAM cells};
\end{scope}
\end{tikzpicture}
\caption{(a) A six-transistor SRAM cell: the word line (WL) turns on the access transistors W1 and W2, connecting the cross-coupled inverters to the complementary bit lines. (b) Cells are arranged in a two-dimensional array; a decoder raises one word line, every cell in that row drives its bit lines, and a column multiplexer and sense amplifiers select and read the bits required. Redrawn after Sarangi \cite[sec.~7.4.1, Figs.~7.26--7.27]{sarangi2025basic}.}
\label{fig:sram-array}
\end{figure}
